%% ECON 282E -- Session 4, Deck B3: Discretizing an AR(1)
%% Discharges S03_A2's promise: "the two standard constructions disagree badly
%% when rho > 0.95 ... We build both next week."
%% Tauchen is ported from QM Lec.4 F69-F75 (with the drift inconsistency fixed --
%% the source gives the transition formula twice, once with mu(1-rho) and once
%% without, and renames the state z -> y -> z).  ROUWENHORST IS NEW: it appears
%% nowhere in the corpus, named twice and never constructed.
%% Numbers from tools/figures/s04_numbers.json, key "markov".

%% ECON 282E -- Foundations of Macroeconomics (UC Riverside, Fall 2026)
%% Shared Beamer preamble for every deck in slides/S01 ... slides/S10.
%%
%% Usage (a deck lives in slides/S0X/ and is compiled from that folder):
%%   %% ECON 282E -- Foundations of Macroeconomics (UC Riverside, Fall 2026)
%% Shared Beamer preamble for every deck in slides/S01 ... slides/S10.
%%
%% Usage (a deck lives in slides/S0X/ and is compiled from that folder):
%%   %% ECON 282E -- Foundations of Macroeconomics (UC Riverside, Fall 2026)
%% Shared Beamer preamble for every deck in slides/S01 ... slides/S10.
%%
%% Usage (a deck lives in slides/S0X/ and is compiled from that folder):
%%   \input{../preamble/course_preamble.tex}
%%   \decktitle{1}{A1}{Who I Am \& Why This Course}{September 24, 2026}
%%   \begin{document}
%%   \makedecktitle
%%   ...
%%
%% Adapted from Courses/AI-ECON-2026/source/shared_preamble.tex (Metropolis theme,
%% concept/caution/example/takeaway boxes, \sectiondivider). Changes for this course:
%%   * course title block (\decktitle / \makedecktitle) instead of \makebeamertitle
%%   * \zh{...} typesets Chinese (book titles) under pdflatex via CJKutf8
%%   * researchbox for the "Research in the wild" frames
%%   * section pages off; TikZ libraries and the course map loaded here
%% Build with pdflatex only (two passes); tools/build_slides.py does this for you.

\documentclass[10pt,english,aspectratio=169]{beamer}

\usepackage[T1]{fontenc}
\usepackage{textcomp}
\usepackage[utf8]{inputenc}
\usepackage{babel}
\usepackage{CJKutf8}

\usepackage{amsmath, amssymb, amsthm, graphicx}
\usepackage{booktabs, multirow, tabularx}
\usepackage{tikz}
\usetikzlibrary{positioning,arrows.meta,shapes,calc,fit,backgrounds}
\usepackage{pgfplots}
\pgfplotsset{compat=1.16}
\usepackage[most]{tcolorbox}
\tcbuselibrary{skins, breakable}
\usepackage{listings}
\usepackage{xcolor}

\lstset{
  basicstyle=\ttfamily\small,
  keywordstyle=\color{blue!80!black}\bfseries,
  commentstyle=\color{green!50!black}\itshape,
  stringstyle=\color{orange!80!black},
  backgroundcolor=\color{gray!8},
  frame=single,
  rulecolor=\color{gray!40},
  breaklines=true,
  showstringspaces=false,
  numbers=none,
}

\usetheme{metropolis}
\usefonttheme{professionalfonts}
\metroset{sectionpage=none}

\definecolor{LightGray}{RGB}{230,230,230}
\definecolor{RedTitle}{RGB}{0,0,200}
\definecolor{VeryLightGray}{RGB}{245,245,245}
\definecolor{DarkText}{RGB}{0,0,0}
\definecolor{UCRBlue}{RGB}{0,45,106}
\definecolor{UCRGold}{RGB}{241,171,0}

% Stage colours of the course arc (used by the course map and elsewhere)
\colorlet{StageTrunk}{blue!12}
\colorlet{StageBranch}{cyan!14}
\colorlet{StageMachine}{gray!18}
\colorlet{StageWall}{red!14}
\colorlet{StageTools}{orange!20}
\colorlet{StageData}{green!16}
\colorlet{StageAgents}{violet!14}

\hypersetup{bookmarks=false,breaklinks=true,pdfborder={0 0 0},colorlinks=true,
  linkcolor=DarkText,urlcolor=RedTitle}

\setbeamercolor{background canvas}{bg=white}
\setbeamercolor{title}{fg=RedTitle}
\setbeamercolor{frametitle}{bg=VeryLightGray,fg=DarkText}
\setbeamercolor{normal text}{fg=DarkText}
\setbeamercolor{alerted text}{fg=RedTitle}
\setbeamersize{text margin left=5mm, text margin right=5mm}
\addtobeamertemplate{frametitle}{}{\vspace{-0.5em}}

\setbeamertemplate{navigation symbols}{}
\setbeamertemplate{footline}{%
  \hspace{5mm}{\usebeamerfont{page number in head/foot}\color{black!45}%
  ECON 282E \,$\cdot$\, \insertshortdeck}\hfill%
  \usebeamercolor[fg]{page number in head/foot}%
  \usebeamerfont{page number in head/foot}%
  \insertframenumber\,/\,\inserttotalframenumber\kern1em\vskip2pt%
}

% ---------------------------------------------------------------------------
% Chinese text under pdflatex (book titles etc.)
\newcommand{\zh}[1]{\begin{CJK*}{UTF8}{gbsn}#1\end{CJK*}}

% ---------------------------------------------------------------------------
% Deck title block
%   \decktitle{session}{deck id}{title}{date}
\newcommand{\insertshortdeck}{}
\newcommand{\decksession}{}
\newcommand{\deckid}{}
\newcommand{\decktitletext}{}
\newcommand{\deckdate}{}
\newcommand{\decktitle}[4]{%
  \renewcommand{\decksession}{#1}%
  \renewcommand{\deckid}{#2}%
  \renewcommand{\decktitletext}{#3}%
  \renewcommand{\deckdate}{#4}%
  \renewcommand{\insertshortdeck}{Session #1 \,$\cdot$\, Deck #1.#2}%
  \title{#3}%
  \author{Zhigang Feng}%
  \date{#4}%
}
\newcommand{\makedecktitle}{%
  \begin{frame}[plain,noframenumbering]
    \begin{tikzpicture}[remember picture,overlay]
      \fill[UCRBlue] (current page.north west) rectangle ([yshift=-0.55cm]current page.north east);
      \fill[UCRGold] ([yshift=-0.55cm]current page.north west) rectangle ([yshift=-0.62cm]current page.north east);
      \node[anchor=west, text=white, font=\small\bfseries] at ([xshift=5mm,yshift=-0.275cm]current page.north west)
        {ECON 282E \,$\cdot$\, Foundations of Macroeconomics};
      \node[anchor=east, text=white, font=\small] at ([xshift=-5mm,yshift=-0.275cm]current page.north east)
        {UC Riverside \,$\cdot$\, Fall 2026};
    \end{tikzpicture}
    \vspace*{1.1cm}
    {\usebeamerfont{title}\color{black!55}\large Session \decksession\ \,$\cdot$\, Deck \decksession.\deckid\par}
    \vspace{0.25cm}
    {\usebeamerfont{title}\color{RedTitle}\LARGE\bfseries \decktitletext\par}
    \vspace{0.7cm}
    {\large Zhigang Feng\par}
    \vspace{0.15cm}
    {\color{black!65}\deckdate\par}
    \vfill
    {\footnotesize\itshape\color{black!55}Build it on paper $\;\to\;$ solve it on the machine $\;\to\;$ push it to the frontier with AI\par}
  \end{frame}}

% ---------------------------------------------------------------------------
% Section divider (plain frame, not counted)
\newcommand{\sectiondivider}[2]{%
\begin{frame}[plain,noframenumbering]
\begin{center}
\vfill
{\Large\textbf{#1}\par}
\vspace{0.35cm}
{\normalsize #2\par}
\vfill
\end{center}
\end{frame}
}

% ---------------------------------------------------------------------------
% Boxes
\newtcolorbox{conceptbox}[1][]{colback=blue!3, colframe=RedTitle!55, boxrule=0.35mm,
  arc=1mm, left=2mm, right=2mm, top=1mm, bottom=1mm, #1}
\newtcolorbox{cautionbox}[1][]{colback=red!3, colframe=red!50!black, boxrule=0.35mm,
  arc=1mm, left=2mm, right=2mm, top=1mm, bottom=1mm, #1}
\newtcolorbox{examplebox}[1][]{colback=green!3, colframe=green!45!black, boxrule=0.35mm,
  arc=1mm, left=2mm, right=2mm, top=1mm, bottom=1mm, #1}
\newtcolorbox{takeawaybox}[1][]{colback=VeryLightGray, colframe=RedTitle!55, boxrule=0.35mm,
  arc=1mm, left=2mm, right=2mm, top=1mm, bottom=1mm, #1}
\newtcolorbox{researchbox}[1][]{enhanced, colback=violet!4, colframe=violet!60!black,
  boxrule=0.35mm, arc=1mm, left=2mm, right=2mm, top=1mm, bottom=1mm,
  fonttitle=\bfseries\small, title={Research in the wild}, #1}

\newcommand{\compacttables}{\renewcommand{\arraystretch}{1.15}}
\newcommand{\src}[1]{{\tiny\color{black!55}#1}}

% ---------------------------------------------------------------------------
\input{../preamble/course_map.tex}

%%   \decktitle{1}{A1}{Who I Am \& Why This Course}{September 24, 2026}
%%   \begin{document}
%%   \makedecktitle
%%   ...
%%
%% Adapted from Courses/AI-ECON-2026/source/shared_preamble.tex (Metropolis theme,
%% concept/caution/example/takeaway boxes, \sectiondivider). Changes for this course:
%%   * course title block (\decktitle / \makedecktitle) instead of \makebeamertitle
%%   * \zh{...} typesets Chinese (book titles) under pdflatex via CJKutf8
%%   * researchbox for the "Research in the wild" frames
%%   * section pages off; TikZ libraries and the course map loaded here
%% Build with pdflatex only (two passes); tools/build_slides.py does this for you.

\documentclass[10pt,english,aspectratio=169]{beamer}

\usepackage[T1]{fontenc}
\usepackage{textcomp}
\usepackage[utf8]{inputenc}
\usepackage{babel}
\usepackage{CJKutf8}

\usepackage{amsmath, amssymb, amsthm, graphicx}
\usepackage{booktabs, multirow, tabularx}
\usepackage{tikz}
\usetikzlibrary{positioning,arrows.meta,shapes,calc,fit,backgrounds}
\usepackage{pgfplots}
\pgfplotsset{compat=1.16}
\usepackage[most]{tcolorbox}
\tcbuselibrary{skins, breakable}
\usepackage{listings}
\usepackage{xcolor}

\lstset{
  basicstyle=\ttfamily\small,
  keywordstyle=\color{blue!80!black}\bfseries,
  commentstyle=\color{green!50!black}\itshape,
  stringstyle=\color{orange!80!black},
  backgroundcolor=\color{gray!8},
  frame=single,
  rulecolor=\color{gray!40},
  breaklines=true,
  showstringspaces=false,
  numbers=none,
}

\usetheme{metropolis}
\usefonttheme{professionalfonts}
\metroset{sectionpage=none}

\definecolor{LightGray}{RGB}{230,230,230}
\definecolor{RedTitle}{RGB}{0,0,200}
\definecolor{VeryLightGray}{RGB}{245,245,245}
\definecolor{DarkText}{RGB}{0,0,0}
\definecolor{UCRBlue}{RGB}{0,45,106}
\definecolor{UCRGold}{RGB}{241,171,0}

% Stage colours of the course arc (used by the course map and elsewhere)
\colorlet{StageTrunk}{blue!12}
\colorlet{StageBranch}{cyan!14}
\colorlet{StageMachine}{gray!18}
\colorlet{StageWall}{red!14}
\colorlet{StageTools}{orange!20}
\colorlet{StageData}{green!16}
\colorlet{StageAgents}{violet!14}

\hypersetup{bookmarks=false,breaklinks=true,pdfborder={0 0 0},colorlinks=true,
  linkcolor=DarkText,urlcolor=RedTitle}

\setbeamercolor{background canvas}{bg=white}
\setbeamercolor{title}{fg=RedTitle}
\setbeamercolor{frametitle}{bg=VeryLightGray,fg=DarkText}
\setbeamercolor{normal text}{fg=DarkText}
\setbeamercolor{alerted text}{fg=RedTitle}
\setbeamersize{text margin left=5mm, text margin right=5mm}
\addtobeamertemplate{frametitle}{}{\vspace{-0.5em}}

\setbeamertemplate{navigation symbols}{}
\setbeamertemplate{footline}{%
  \hspace{5mm}{\usebeamerfont{page number in head/foot}\color{black!45}%
  ECON 282E \,$\cdot$\, \insertshortdeck}\hfill%
  \usebeamercolor[fg]{page number in head/foot}%
  \usebeamerfont{page number in head/foot}%
  \insertframenumber\,/\,\inserttotalframenumber\kern1em\vskip2pt%
}

% ---------------------------------------------------------------------------
% Chinese text under pdflatex (book titles etc.)
\newcommand{\zh}[1]{\begin{CJK*}{UTF8}{gbsn}#1\end{CJK*}}

% ---------------------------------------------------------------------------
% Deck title block
%   \decktitle{session}{deck id}{title}{date}
\newcommand{\insertshortdeck}{}
\newcommand{\decksession}{}
\newcommand{\deckid}{}
\newcommand{\decktitletext}{}
\newcommand{\deckdate}{}
\newcommand{\decktitle}[4]{%
  \renewcommand{\decksession}{#1}%
  \renewcommand{\deckid}{#2}%
  \renewcommand{\decktitletext}{#3}%
  \renewcommand{\deckdate}{#4}%
  \renewcommand{\insertshortdeck}{Session #1 \,$\cdot$\, Deck #1.#2}%
  \title{#3}%
  \author{Zhigang Feng}%
  \date{#4}%
}
\newcommand{\makedecktitle}{%
  \begin{frame}[plain,noframenumbering]
    \begin{tikzpicture}[remember picture,overlay]
      \fill[UCRBlue] (current page.north west) rectangle ([yshift=-0.55cm]current page.north east);
      \fill[UCRGold] ([yshift=-0.55cm]current page.north west) rectangle ([yshift=-0.62cm]current page.north east);
      \node[anchor=west, text=white, font=\small\bfseries] at ([xshift=5mm,yshift=-0.275cm]current page.north west)
        {ECON 282E \,$\cdot$\, Foundations of Macroeconomics};
      \node[anchor=east, text=white, font=\small] at ([xshift=-5mm,yshift=-0.275cm]current page.north east)
        {UC Riverside \,$\cdot$\, Fall 2026};
    \end{tikzpicture}
    \vspace*{1.1cm}
    {\usebeamerfont{title}\color{black!55}\large Session \decksession\ \,$\cdot$\, Deck \decksession.\deckid\par}
    \vspace{0.25cm}
    {\usebeamerfont{title}\color{RedTitle}\LARGE\bfseries \decktitletext\par}
    \vspace{0.7cm}
    {\large Zhigang Feng\par}
    \vspace{0.15cm}
    {\color{black!65}\deckdate\par}
    \vfill
    {\footnotesize\itshape\color{black!55}Build it on paper $\;\to\;$ solve it on the machine $\;\to\;$ push it to the frontier with AI\par}
  \end{frame}}

% ---------------------------------------------------------------------------
% Section divider (plain frame, not counted)
\newcommand{\sectiondivider}[2]{%
\begin{frame}[plain,noframenumbering]
\begin{center}
\vfill
{\Large\textbf{#1}\par}
\vspace{0.35cm}
{\normalsize #2\par}
\vfill
\end{center}
\end{frame}
}

% ---------------------------------------------------------------------------
% Boxes
\newtcolorbox{conceptbox}[1][]{colback=blue!3, colframe=RedTitle!55, boxrule=0.35mm,
  arc=1mm, left=2mm, right=2mm, top=1mm, bottom=1mm, #1}
\newtcolorbox{cautionbox}[1][]{colback=red!3, colframe=red!50!black, boxrule=0.35mm,
  arc=1mm, left=2mm, right=2mm, top=1mm, bottom=1mm, #1}
\newtcolorbox{examplebox}[1][]{colback=green!3, colframe=green!45!black, boxrule=0.35mm,
  arc=1mm, left=2mm, right=2mm, top=1mm, bottom=1mm, #1}
\newtcolorbox{takeawaybox}[1][]{colback=VeryLightGray, colframe=RedTitle!55, boxrule=0.35mm,
  arc=1mm, left=2mm, right=2mm, top=1mm, bottom=1mm, #1}
\newtcolorbox{researchbox}[1][]{enhanced, colback=violet!4, colframe=violet!60!black,
  boxrule=0.35mm, arc=1mm, left=2mm, right=2mm, top=1mm, bottom=1mm,
  fonttitle=\bfseries\small, title={Research in the wild}, #1}

\newcommand{\compacttables}{\renewcommand{\arraystretch}{1.15}}
\newcommand{\src}[1]{{\tiny\color{black!55}#1}}

% ---------------------------------------------------------------------------
%% Course map: the recurring "you are here" slide for ECON 282E.
%% Loaded by course_preamble.tex. Pattern follows AI-ECON-2026 lec01_map.tex, but the map
%% shows the whole ten-session course rather than one lecture.
%%
%%   \CourseMap{s}{label}   s = 1..10 highlights session s and prints "you are here: label"
%%                          (e.g. \CourseMap{1}{1.A2}); earlier sessions get a check mark.
%%   \CourseMap{0}{}        full overview, nothing highlighted.
%%
%% Note: TikZ option lists are comma-split before TeX conditionals run, so every \ifnum
%% below sits outside [...] option lists.

\newcommand{\CourseMap}[2]{%
\begin{center}
\begin{tikzpicture}[
    sess/.style={rectangle, rounded corners=2pt, draw=black!45, minimum width=1.34cm,
        minimum height=1.12cm, text width=1.26cm, align=center, inner sep=1pt},
    stage/.style={font=\tiny\bfseries, text=black!70},
]
  \foreach \i/\d/\t/\c in {%
      1/Sep 24/Intro \\ Growth/StageTrunk,
      2/Oct 1/HA $\cdot$ OLG \\ Policy/StageBranch,
      3/Oct 8/Num.\ DP \\ Python/StageMachine,
      4/Oct 15/Numerics \\ I/StageMachine,
      5/Oct 22/Numerics \\ II/StageMachine,
      6/Oct 29/The wall \\ \textbf{Midterm}/StageWall,
      7/Nov 5/ML \\ Deep nets/StageTools,
      8/Nov 12/RL \\ HA $+$ DL/StageTools,
      9/Nov 19/Text \\ LLMs/StageData,
      10/Dec 3/Agents \\ \textbf{Projects}/StageAgents} {
    \node[sess, fill=\c] (n\i) at ({1.46*(\i-1)}, 0)
      {{\scriptsize\bfseries S\i}\\[-1pt]{\tiny\color{black!60}\d}\\[1pt]{\tiny \t}};
    \ifnum\i<#1
      \node[font=\scriptsize, text=green!45!black] at ([xshift=-2.5mm,yshift=-2.2mm]n\i.north east) {$\checkmark$};
    \fi
    \ifnum\i=#1
      \draw[RedTitle, very thick, rounded corners=2pt]
        ([xshift=-0.6mm,yshift=-0.6mm]n\i.south west) rectangle ([xshift=0.6mm,yshift=0.6mm]n\i.north east);
      % keep the label inside the picture at both ends of the row
      \ifnum\i<3
        \node[font=\scriptsize\bfseries, text=RedTitle, anchor=south west] at ([yshift=1.2mm]n\i.north west)
          {$\blacktriangledown$\,you are here: #2};
      \else\ifnum\i>8
        \node[font=\scriptsize\bfseries, text=RedTitle, anchor=south east] at ([yshift=1.2mm]n\i.north east)
          {you are here: #2\,$\blacktriangledown$};
      \else
        \node[font=\scriptsize\bfseries, text=RedTitle, anchor=south] at ([yshift=1.2mm]n\i.north)
          {you are here: #2\,$\blacktriangledown$};
      \fi\fi
    \fi
  }
  % stage band under the sessions
  \foreach \a/\b/\lab/\c in {1/1/trunk/StageTrunk, 2/2/branches/StageBranch,
      3/5/the machine $\cdot$ classical methods/StageMachine, 6/6/the wall/StageWall,
      7/8/new tools/StageTools, 9/9/new data/StageData, 10/10/colleagues/StageAgents} {
    \fill[\c] ([yshift=-1.5mm]n\a.south west) rectangle ([yshift=-4.6mm]n\b.south east);
    \path ([yshift=-1.5mm]n\a.south west) -- ([yshift=-4.6mm]n\b.south east)
      node[midway, stage] {\lab};
  }
  \node[font=\footnotesize\itshape, text=black!70] at ([yshift=-9.5mm]$(n1.south)!0.5!(n10.south)$)
    {Build it on paper $\;\to\;$ solve it on the machine $\;\to\;$ push it to the frontier with AI};
\end{tikzpicture}
\end{center}
}


%%   \decktitle{1}{A1}{Who I Am \& Why This Course}{September 24, 2026}
%%   \begin{document}
%%   \makedecktitle
%%   ...
%%
%% Adapted from Courses/AI-ECON-2026/source/shared_preamble.tex (Metropolis theme,
%% concept/caution/example/takeaway boxes, \sectiondivider). Changes for this course:
%%   * course title block (\decktitle / \makedecktitle) instead of \makebeamertitle
%%   * \zh{...} typesets Chinese (book titles) under pdflatex via CJKutf8
%%   * researchbox for the "Research in the wild" frames
%%   * section pages off; TikZ libraries and the course map loaded here
%% Build with pdflatex only (two passes); tools/build_slides.py does this for you.

\documentclass[10pt,english,aspectratio=169]{beamer}

\usepackage[T1]{fontenc}
\usepackage{textcomp}
\usepackage[utf8]{inputenc}
\usepackage{babel}
\usepackage{CJKutf8}

\usepackage{amsmath, amssymb, amsthm, graphicx}
\usepackage{booktabs, multirow, tabularx}
\usepackage{tikz}
\usetikzlibrary{positioning,arrows.meta,shapes,calc,fit,backgrounds}
\usepackage{pgfplots}
\pgfplotsset{compat=1.16}
\usepackage[most]{tcolorbox}
\tcbuselibrary{skins, breakable}
\usepackage{listings}
\usepackage{xcolor}

\lstset{
  basicstyle=\ttfamily\small,
  keywordstyle=\color{blue!80!black}\bfseries,
  commentstyle=\color{green!50!black}\itshape,
  stringstyle=\color{orange!80!black},
  backgroundcolor=\color{gray!8},
  frame=single,
  rulecolor=\color{gray!40},
  breaklines=true,
  showstringspaces=false,
  numbers=none,
}

\usetheme{metropolis}
\usefonttheme{professionalfonts}
\metroset{sectionpage=none}

\definecolor{LightGray}{RGB}{230,230,230}
\definecolor{RedTitle}{RGB}{0,0,200}
\definecolor{VeryLightGray}{RGB}{245,245,245}
\definecolor{DarkText}{RGB}{0,0,0}
\definecolor{UCRBlue}{RGB}{0,45,106}
\definecolor{UCRGold}{RGB}{241,171,0}

% Stage colours of the course arc (used by the course map and elsewhere)
\colorlet{StageTrunk}{blue!12}
\colorlet{StageBranch}{cyan!14}
\colorlet{StageMachine}{gray!18}
\colorlet{StageWall}{red!14}
\colorlet{StageTools}{orange!20}
\colorlet{StageData}{green!16}
\colorlet{StageAgents}{violet!14}

\hypersetup{bookmarks=false,breaklinks=true,pdfborder={0 0 0},colorlinks=true,
  linkcolor=DarkText,urlcolor=RedTitle}

\setbeamercolor{background canvas}{bg=white}
\setbeamercolor{title}{fg=RedTitle}
\setbeamercolor{frametitle}{bg=VeryLightGray,fg=DarkText}
\setbeamercolor{normal text}{fg=DarkText}
\setbeamercolor{alerted text}{fg=RedTitle}
\setbeamersize{text margin left=5mm, text margin right=5mm}
\addtobeamertemplate{frametitle}{}{\vspace{-0.5em}}

\setbeamertemplate{navigation symbols}{}
\setbeamertemplate{footline}{%
  \hspace{5mm}{\usebeamerfont{page number in head/foot}\color{black!45}%
  ECON 282E \,$\cdot$\, \insertshortdeck}\hfill%
  \usebeamercolor[fg]{page number in head/foot}%
  \usebeamerfont{page number in head/foot}%
  \insertframenumber\,/\,\inserttotalframenumber\kern1em\vskip2pt%
}

% ---------------------------------------------------------------------------
% Chinese text under pdflatex (book titles etc.)
\newcommand{\zh}[1]{\begin{CJK*}{UTF8}{gbsn}#1\end{CJK*}}

% ---------------------------------------------------------------------------
% Deck title block
%   \decktitle{session}{deck id}{title}{date}
\newcommand{\insertshortdeck}{}
\newcommand{\decksession}{}
\newcommand{\deckid}{}
\newcommand{\decktitletext}{}
\newcommand{\deckdate}{}
\newcommand{\decktitle}[4]{%
  \renewcommand{\decksession}{#1}%
  \renewcommand{\deckid}{#2}%
  \renewcommand{\decktitletext}{#3}%
  \renewcommand{\deckdate}{#4}%
  \renewcommand{\insertshortdeck}{Session #1 \,$\cdot$\, Deck #1.#2}%
  \title{#3}%
  \author{Zhigang Feng}%
  \date{#4}%
}
\newcommand{\makedecktitle}{%
  \begin{frame}[plain,noframenumbering]
    \begin{tikzpicture}[remember picture,overlay]
      \fill[UCRBlue] (current page.north west) rectangle ([yshift=-0.55cm]current page.north east);
      \fill[UCRGold] ([yshift=-0.55cm]current page.north west) rectangle ([yshift=-0.62cm]current page.north east);
      \node[anchor=west, text=white, font=\small\bfseries] at ([xshift=5mm,yshift=-0.275cm]current page.north west)
        {ECON 282E \,$\cdot$\, Foundations of Macroeconomics};
      \node[anchor=east, text=white, font=\small] at ([xshift=-5mm,yshift=-0.275cm]current page.north east)
        {UC Riverside \,$\cdot$\, Fall 2026};
    \end{tikzpicture}
    \vspace*{1.1cm}
    {\usebeamerfont{title}\color{black!55}\large Session \decksession\ \,$\cdot$\, Deck \decksession.\deckid\par}
    \vspace{0.25cm}
    {\usebeamerfont{title}\color{RedTitle}\LARGE\bfseries \decktitletext\par}
    \vspace{0.7cm}
    {\large Zhigang Feng\par}
    \vspace{0.15cm}
    {\color{black!65}\deckdate\par}
    \vfill
    {\footnotesize\itshape\color{black!55}Build it on paper $\;\to\;$ solve it on the machine $\;\to\;$ push it to the frontier with AI\par}
  \end{frame}}

% ---------------------------------------------------------------------------
% Section divider (plain frame, not counted)
\newcommand{\sectiondivider}[2]{%
\begin{frame}[plain,noframenumbering]
\begin{center}
\vfill
{\Large\textbf{#1}\par}
\vspace{0.35cm}
{\normalsize #2\par}
\vfill
\end{center}
\end{frame}
}

% ---------------------------------------------------------------------------
% Boxes
\newtcolorbox{conceptbox}[1][]{colback=blue!3, colframe=RedTitle!55, boxrule=0.35mm,
  arc=1mm, left=2mm, right=2mm, top=1mm, bottom=1mm, #1}
\newtcolorbox{cautionbox}[1][]{colback=red!3, colframe=red!50!black, boxrule=0.35mm,
  arc=1mm, left=2mm, right=2mm, top=1mm, bottom=1mm, #1}
\newtcolorbox{examplebox}[1][]{colback=green!3, colframe=green!45!black, boxrule=0.35mm,
  arc=1mm, left=2mm, right=2mm, top=1mm, bottom=1mm, #1}
\newtcolorbox{takeawaybox}[1][]{colback=VeryLightGray, colframe=RedTitle!55, boxrule=0.35mm,
  arc=1mm, left=2mm, right=2mm, top=1mm, bottom=1mm, #1}
\newtcolorbox{researchbox}[1][]{enhanced, colback=violet!4, colframe=violet!60!black,
  boxrule=0.35mm, arc=1mm, left=2mm, right=2mm, top=1mm, bottom=1mm,
  fonttitle=\bfseries\small, title={Research in the wild}, #1}

\newcommand{\compacttables}{\renewcommand{\arraystretch}{1.15}}
\newcommand{\src}[1]{{\tiny\color{black!55}#1}}

% ---------------------------------------------------------------------------
%% Course map: the recurring "you are here" slide for ECON 282E.
%% Loaded by course_preamble.tex. Pattern follows AI-ECON-2026 lec01_map.tex, but the map
%% shows the whole ten-session course rather than one lecture.
%%
%%   \CourseMap{s}{label}   s = 1..10 highlights session s and prints "you are here: label"
%%                          (e.g. \CourseMap{1}{1.A2}); earlier sessions get a check mark.
%%   \CourseMap{0}{}        full overview, nothing highlighted.
%%
%% Note: TikZ option lists are comma-split before TeX conditionals run, so every \ifnum
%% below sits outside [...] option lists.

\newcommand{\CourseMap}[2]{%
\begin{center}
\begin{tikzpicture}[
    sess/.style={rectangle, rounded corners=2pt, draw=black!45, minimum width=1.34cm,
        minimum height=1.12cm, text width=1.26cm, align=center, inner sep=1pt},
    stage/.style={font=\tiny\bfseries, text=black!70},
]
  \foreach \i/\d/\t/\c in {%
      1/Sep 24/Intro \\ Growth/StageTrunk,
      2/Oct 1/HA $\cdot$ OLG \\ Policy/StageBranch,
      3/Oct 8/Num.\ DP \\ Python/StageMachine,
      4/Oct 15/Numerics \\ I/StageMachine,
      5/Oct 22/Numerics \\ II/StageMachine,
      6/Oct 29/The wall \\ \textbf{Midterm}/StageWall,
      7/Nov 5/ML \\ Deep nets/StageTools,
      8/Nov 12/RL \\ HA $+$ DL/StageTools,
      9/Nov 19/Text \\ LLMs/StageData,
      10/Dec 3/Agents \\ \textbf{Projects}/StageAgents} {
    \node[sess, fill=\c] (n\i) at ({1.46*(\i-1)}, 0)
      {{\scriptsize\bfseries S\i}\\[-1pt]{\tiny\color{black!60}\d}\\[1pt]{\tiny \t}};
    \ifnum\i<#1
      \node[font=\scriptsize, text=green!45!black] at ([xshift=-2.5mm,yshift=-2.2mm]n\i.north east) {$\checkmark$};
    \fi
    \ifnum\i=#1
      \draw[RedTitle, very thick, rounded corners=2pt]
        ([xshift=-0.6mm,yshift=-0.6mm]n\i.south west) rectangle ([xshift=0.6mm,yshift=0.6mm]n\i.north east);
      % keep the label inside the picture at both ends of the row
      \ifnum\i<3
        \node[font=\scriptsize\bfseries, text=RedTitle, anchor=south west] at ([yshift=1.2mm]n\i.north west)
          {$\blacktriangledown$\,you are here: #2};
      \else\ifnum\i>8
        \node[font=\scriptsize\bfseries, text=RedTitle, anchor=south east] at ([yshift=1.2mm]n\i.north east)
          {you are here: #2\,$\blacktriangledown$};
      \else
        \node[font=\scriptsize\bfseries, text=RedTitle, anchor=south] at ([yshift=1.2mm]n\i.north)
          {you are here: #2\,$\blacktriangledown$};
      \fi\fi
    \fi
  }
  % stage band under the sessions
  \foreach \a/\b/\lab/\c in {1/1/trunk/StageTrunk, 2/2/branches/StageBranch,
      3/5/the machine $\cdot$ classical methods/StageMachine, 6/6/the wall/StageWall,
      7/8/new tools/StageTools, 9/9/new data/StageData, 10/10/colleagues/StageAgents} {
    \fill[\c] ([yshift=-1.5mm]n\a.south west) rectangle ([yshift=-4.6mm]n\b.south east);
    \path ([yshift=-1.5mm]n\a.south west) -- ([yshift=-4.6mm]n\b.south east)
      node[midway, stage] {\lab};
  }
  \node[font=\footnotesize\itshape, text=black!70] at ([yshift=-9.5mm]$(n1.south)!0.5!(n10.south)$)
    {Build it on paper $\;\to\;$ solve it on the machine $\;\to\;$ push it to the frontier with AI};
\end{tikzpicture}
\end{center}
}


\graphicspath{{./figures/}}

\lstset{language=Python, basicstyle=\ttfamily\scriptsize, frame=none,
        backgroundcolor=\color{gray!8}, aboveskip=2pt, belowskip=2pt}

\decktitle{4}{B3}{Discretizing an AR(1): Where $\Pi$ Comes From}{Thursday, October 15, 2026}

\begin{document}

\makedecktitle

%=============================================================================
\begin{frame}{Where We Are}
\CourseMap{4}{4.B3}
\vspace{-1mm}
\begin{center}
\small \textbf{Block B:} \textbf{4.B1} approximation and interpolation $\;\cdot\;$ \textbf{4.B2} quadrature $\;\cdot\;$ \textbf{4.B3} discretizing an AR(1) $\;\cdot\;$ \textbf{4.B4} fast and accurate $\;\cdot\;$ \textbf{4.B5} the worked example, end to end.
\end{center}
\end{frame}

%=============================================================================
\begin{frame}{A Chain That Reports a Unit Root}
\begin{columns}[T]
\begin{column}{0.51\textwidth}
\small
Take the standard quarterly productivity process,
\[
  \ln z'=\rho\ln z+\sigma\varepsilon', \qquad \rho=0.99,\;\sigma=0.007,
\]
and build a five-state chain the usual way. Then ask the chain what persistence it actually has.
\medskip
\begin{center}\footnotesize
\begin{tabularx}{\linewidth}{@{}>{\raggedright\arraybackslash}X c c@{}}
\toprule
$\rho=0.99$, $n=5$ & \textbf{implied $\rho$} & \textbf{sd error} \\
\midrule
Tauchen & \textbf{1.0000} & $36\%$ \\
Rouwenhorst & 0.990000 & $6\times10^{-13}\%$ \\
\bottomrule
\end{tabularx}
\end{center}
\medskip
The Tauchen chain is a \textbf{unit root}, and its unconditional standard deviation is a third too small. Every business-cycle moment computed from it is wrong.
\end{column}
\begin{column}{0.46\textwidth}
\begin{cautionbox}[title=\textbf{Session 3 owed you this}]\small
3.A2 said the two standard constructions ``disagree badly when $\rho>0.95$, which is exactly the range macro uses''. This is the measurement.
\end{cautionbox}
\medskip
\begin{takeawaybox}\footnotesize
And $\rho=0.99$ is not an exotic calibration. It is the \textbf{standard quarterly} one --- the same $\rho$ used in this session's RBC model and in Homework 1.
\end{takeawaybox}
\end{column}
\end{columns}
\end{frame}

%=============================================================================
\begin{frame}{What We Are Actually Asking For}
\begin{columns}[T]
\begin{column}{0.53\textwidth}
\small
A continuous AR(1) has a continuum of states. The Bellman equation needs a finite sum, so we replace it by a chain: a grid $\{z_1,\dots,z_n\}$ and a matrix $\Pi$ with
\[
  \Pi_{jm}=\Pr\big(z'=z_m \mid z=z_j\big).
\]
\medskip
This is \emph{not} the quadrature problem of 4.B2. There the distribution was fixed; here it \textbf{depends on today's state}, and the object we need is a whole conditional law for each $j$.
\medskip
\textbf{What a good chain should match:} the persistence $\rho$, the unconditional variance $\sigma^{2}/(1-\rho^{2})$, and ideally the shape of the conditional distribution.
\end{column}
\begin{column}{0.44\textwidth}
\begin{conceptbox}[title=\textbf{Why $\rho$ near one is hard}]\small
The unconditional standard deviation is
\[
  \sigma_z=\frac{\sigma}{\sqrt{1-\rho^{2}}},
\]
which \textbf{explodes} as $\rho\to1$: at $\rho=0.99$ it is seven times $\sigma$. A fixed number of states has to cover a much wider range, with each step correspondingly coarser.
\end{conceptbox}
\end{column}
\end{columns}
\vspace{1mm}
\src{Setup adapted from QM Lec.~4, ``AR(1) Discretization: Why?'' and ``AR(1) Process: Properties''.}
\end{frame}

%=============================================================================
\begin{frame}{Tauchen (1986): Put a Grid Down and Integrate the Normal}
\begin{columns}[T]
\begin{column}{0.53\textwidth}
\small
\begin{enumerate}\footnotesize
  \item Set $\sigma_z=\sigma/\sqrt{1-\rho^{2}}$ and cover $[-m\sigma_z,\,m\sigma_z]$ with $n$ \textbf{equally spaced} points, $m=3$ by convention.
  \item Conditional on $z_j$, next period is $N(\rho z_j,\ \sigma^{2})$.
  \item Assign to each interior $z_m$ the probability that the normal lands in its cell:
\end{enumerate}
\[
  \Pi_{jm}=\Phi\!\Big(\tfrac{z_m+\frac{\Delta}{2}-\rho z_j}{\sigma}\Big)
          -\Phi\!\Big(\tfrac{z_m-\frac{\Delta}{2}-\rho z_j}{\sigma}\Big),
\]
with the two end states absorbing the remaining tails.
\end{column}
\begin{column}{0.44\textwidth}
\begin{cautionbox}[title=\textbf{Two things to fix on sight}]\footnotesize
The formula is often printed twice in the same set of notes, once with a drift term $\mu(1-\rho)$ and once without, with the state renamed halfway through. \textbf{Set $\mu=0$ and keep one symbol.}\\[1mm]
And the coverage multiple is usually called $\lambda$ --- which this course reserves for the invariant distribution. We write $m$.
\end{cautionbox}
\medskip
\begin{examplebox}\footnotesize
The method is intuitive and it is what almost everyone writes first. Its weakness is structural: the grid is chosen \emph{before} the persistence is taken into account.
\end{examplebox}
\end{column}
\end{columns}
\vspace{1mm}
\src{Tauchen (1986). Algorithm from QM Lec.~4, ``Tauchen Method'' frames, with the drift inconsistency resolved.}
\end{frame}

%=============================================================================
\begin{frame}{Rouwenhorst (1995): Build the Chain, Not the Integral}
\begin{columns}[T]
\begin{column}{0.52\textwidth}
\small
A completely different idea: never integrate anything. Start from the two-state chain
\[
  \Pi_2=\begin{pmatrix}p & 1-p\\ 1-p & p\end{pmatrix},
  \qquad p=\frac{1+\rho}{2},
\]
and grow it recursively. Given $\Pi_{n-1}$,
\[
  \Pi_n=p\begin{pmatrix}\Pi_{n-1}&0\\0&0\end{pmatrix}
       +(1-p)\begin{pmatrix}0&\Pi_{n-1}\\0&0\end{pmatrix}
\]
\[
  +(1-p)\begin{pmatrix}0&0\\ \Pi_{n-1}&0\end{pmatrix}
       +p\begin{pmatrix}0&0\\0&\Pi_{n-1}\end{pmatrix},
\]
then halve the interior rows so they sum to one. The grid is equally spaced on $[-\psi,\psi]$ with $\psi=\sigma_z\sqrt{n-1}$.
\end{column}
\begin{column}{0.45\textwidth}
\begin{takeawaybox}[title=\textbf{Why it is exact}]\small
The construction is a \textbf{binomial} one: the chain is a sum of $n-1$ independent two-state processes. Its persistence is $2p-1=\rho$ and its variance is $\psi^{2}/(n-1)=\sigma_z^{2}$ --- \emph{by construction}, for every $n$ and every $\rho$.
\end{takeawaybox}
\medskip
\begin{cautionbox}\footnotesize
The price is that the conditional distribution is binomial, not normal. For a linear-quadratic world that is irrelevant; for a strongly non-linear one it can matter.
\end{cautionbox}
\end{column}
\end{columns}
\vspace{1mm}
\src{Rouwenhorst (1995). \textbf{Written for this course} --- the method is named twice in the source decks and never constructed.}
\end{frame}

%=============================================================================
\begin{frame}{The Comparison Session 3 Promised}
\begin{columns}[c]
\begin{column}{0.53\textwidth}
\centering
\begin{tikzpicture}
\begin{axis}[
    width=7.9cm, height=5.5cm, xmin=3, xmax=17, ymin=-0.3, ymax=4.6,
    xlabel={\scriptsize number of states $n$},
    ylabel={\scriptsize error in implied $\rho$ (\%)},
    tick label style={font=\scriptsize}, xtick={5,9,15},
    axis x line*=bottom, axis y line*=left,
    legend style={font=\tiny, draw=none, fill=none, at={(0.98,0.97)}, anchor=north east}]
  \addplot[very thick, RedTitle, mark=*, mark size=1.3pt] coordinates {(5,3.5028) (9,0.1757) (15,0.1470)};
  \addlegendentry{Tauchen, $\rho=0.90$}
  \addplot[very thick, orange!80!black, mark=square*, mark size=1.2pt] coordinates {(5,3.9866) (9,0.1347) (15,0.1169)};
  \addlegendentry{Tauchen, $\rho=0.95$}
  \addplot[very thick, violet!70!black, mark=triangle*, mark size=1.4pt] coordinates {(5,1.0101) (9,0.8718) (15,0.1767)};
  \addlegendentry{Tauchen, $\rho=0.99$}
  \addplot[very thick, black!55, dashed] coordinates {(5,0) (9,0) (15,0)};
  \addlegendentry{Rouwenhorst (all $\rho$)}
\end{axis}
\end{tikzpicture}
\end{column}
\begin{column}{0.44\textwidth}
\small
Error in the chain's implied persistence, as a percentage of the target $\rho$.
\medskip
\begin{center}\scriptsize
\begin{tabularx}{\linewidth}{@{}c c r r@{}}
\toprule
$\rho$ & $n$ & \textbf{Tauchen} & \textbf{Rouwen.} \\
\midrule
0.90 & 5 & 3.50\% & $2\!\times\!10^{-14}$ \\
0.95 & 5 & 3.99\% & $0$ \\
0.99 & 5 & 1.01\% & $0$ \\
0.99 & 15 & 0.18\% & $7\!\times\!10^{-14}$ \\
\bottomrule
\end{tabularx}
\end{center}
\vspace{1mm}
\begin{cautionbox}\footnotesize
Rouwenhorst is exact to machine precision at \textbf{every} $n$ and \textbf{every} $\rho$. The two lines are not close; one of them is zero.
\end{cautionbox}
\end{column}
\end{columns}
\end{frame}

%=============================================================================
\begin{frame}{And the Variance Is Worse Than the Persistence}
\begin{columns}[T]
\begin{column}{0.53\textwidth}
\small
Persistence is the forgiving statistic. The unconditional standard deviation --- which is what sets the \emph{size} of every business-cycle moment --- is where Tauchen really suffers.
\medskip
\begin{center}\footnotesize
\begin{tabularx}{\linewidth}{@{}c c r r@{}}
\toprule
\multicolumn{2}{@{}l}{sd error} & \textbf{Tauchen} & \textbf{Rouwenhorst} \\
\midrule
$\rho=0.90$ & $n=5$ & 26.9\% & $\approx0$ \\
$\rho=0.95$ & $n=5$ & 31.4\% & $\approx0$ \\
$\rho=0.99$ & $n=5$ & 36.2\% & $\approx0$ \\
$\rho=0.99$ & $n=9$ & 28.5\% & $\approx0$ \\
$\rho=0.99$ & $n=15$ & \textbf{20.1\%} & $\approx0$ \\
\bottomrule
\end{tabularx}
\end{center}
\end{column}
\begin{column}{0.44\textwidth}
\begin{cautionbox}[title=\textbf{Read the last row}]\small
Even with \textbf{fifteen} states, at the standard quarterly persistence, the Tauchen chain's unconditional standard deviation is \textbf{a fifth too small}.
\end{cautionbox}
\medskip
\begin{takeawaybox}\footnotesize
A model calibrated with that chain will understate the volatility of output --- and the natural response, raising $\sigma$ until the moments match, hides a discretization error inside a structural parameter.
\end{takeawaybox}
\end{column}
\end{columns}
\vspace{1mm}
\src{Target $\sigma_z=\sigma/\sqrt{1-\rho^{2}}$; chain moments computed from the invariant distribution of each $\Pi$.}
\end{frame}

%=============================================================================
\begin{frame}{What To Actually Do}
\begin{columns}[T]
\begin{column}{0.51\textwidth}
\renewcommand{\arraystretch}{1.15}
\begin{center}\footnotesize
\begin{tabularx}{\linewidth}{@{}>{\raggedright\arraybackslash}p{2.2cm} >{\raggedright\arraybackslash}X@{}}
\toprule
\textbf{Method} & \textbf{Use it when} \\
\midrule
Rouwenhorst & \textbf{the default}, and essential for $\rho>0.9$ \\
Tauchen & low persistence, or when the conditional distribution must look normal \\
Tauchen--Hussey & quadrature-based; accurate for moderate $\rho$, degrades badly near one \\
\bottomrule
\end{tabularx}
\end{center}
\medskip
\footnotesize
Kopecky \& Suen (2010) make the case formally: Rouwenhorst matches the conditional and unconditional first and second moments and the persistence \emph{exactly}, for any $n$.
\end{column}
\begin{column}{0.46\textwidth}
\begin{takeawaybox}[title=\textbf{The habit to build}]\small
Whatever you use, \textbf{report what the chain actually delivers}: its implied $\rho$ and $\sigma_z$ next to the targets. It is three lines of code and it is the only way anyone can tell whether your moments came from the economics or the grid.
\end{takeawaybox}
\medskip
\begin{examplebox}\footnotesize
Homework 1 Part B asks for exactly that table at $n=5,9,15$ --- and then asks you to solve the model with both and see whether it moved the answer.
\end{examplebox}
\end{column}
\end{columns}
\vspace{1mm}
\src{Kopecky \& Suen, ``Finite State Markov-Chain Approximations to Highly Persistent Processes'', \emph{RED} (2010).}
\end{frame}

%=============================================================================
\begin{frame}{Takeaways}
\begin{takeawaybox}
\begin{itemize}\small
  \item $\Pi$ is not data and it is not a modelling choice you can leave implicit. It is an \textbf{approximation with an error}, and the error is measurable.
  \item \textbf{Tauchen} puts down an equally spaced grid and integrates the conditional normal over each cell. Intuitive, and the grid is chosen before persistence is taken into account.
  \item \textbf{Rouwenhorst} builds the chain recursively from a two-state block and matches $\rho$ and $\sigma_z$ \textbf{exactly, by construction}, for every $n$.
  \item Measured at the standard quarterly calibration $\rho=0.99$: the five-state Tauchen chain implies $\rho=1.0000$ --- \textbf{a unit root} --- and a standard deviation 36\% too small. At $n=15$ it is still 20\% too small.
  \item \textbf{Use Rouwenhorst}, and report the chain's implied $\rho$ and $\sigma_z$ beside the targets every time.
\end{itemize}
\end{takeawaybox}
\end{frame}

%=============================================================================
\begin{frame}{Bridge: Everything Is Now in Place}
\begin{columns}[T]
\begin{column}{0.53\textwidth}
\small
Session 3 ended with six crude choices and a list of what would repair each. We have now built all of them: a root-finder, a maximizer on an interval, derivatives, an interpolation scheme that keeps its shape, a quadrature rule, and a transition matrix that is not invented.
\medskip

What we have not done is put them back together. The crude solver of 3.A4 took 1{,}352 sweeps and 2.4 billion comparisons at this session's calibration, and its worst Euler error was $10^{-1.5}$. The question for the last deck is simple: \textbf{how much of that was avoidable, and which repair bought the most?}
\end{column}
\begin{column}{0.44\textwidth}
\begin{conceptbox}[title=\textbf{Next (4.B4)}]\small
Howard's policy improvement, monotonicity, concavity --- and the endogenous grid method, done the way Carroll actually formulated it, which removes the root-finding entirely.\\[1mm]
Then one table: every technique, its speed-up and its accuracy, measured on the same model.
\end{conceptbox}
\end{column}
\end{columns}
\end{frame}

\end{document}
